\section{Degrees on curves}
\label{section-divisors-curves}

\noindent
We start defining the degree of an invertible sheaf and more generally
a locally free sheaf on a proper scheme
of dimension $1$ over a field. In Section \ref{section-euler}
we defined the Euler characteristic
of a coherent sheaf $\mathcal{F}$ on a proper scheme $X$ over a field
$k$ by the formula
$$
\chi(X, \mathcal{F}) = \sum (-1)^i \dim_k H^i(X, \mathcal{F}).
$$

\begin{definition}
\label{definition-degree-invertible-sheaf}
Let $k$ be a field, let $X$ be a proper scheme of dimension $\leq 1$
over $k$, and let $\mathcal{L}$ be an invertible $\mathcal{O}_X$-module.
The {\it degree} of $\mathcal{L}$ is defined by
$$
\deg(\mathcal{L}) = \chi(X, \mathcal{L}) - \chi(X, \mathcal{O}_X)
$$
More generally, if $\mathcal{E}$ is a locally free sheaf of rank $n$
we define the {\it degree} of $\mathcal{E}$ by
$$
\deg(\mathcal{E}) = \chi(X, \mathcal{E}) - n\chi(X, \mathcal{O}_X)
$$
\end{definition}

\noindent
Observe that this depends on the triple $\mathcal{E}/X/k$.
If $X$ is disconnected and $\mathcal{E}$ is finite locally free (but not
of constant rank), then one can modify the definition by summing the degrees
of the restriction of $\mathcal{E}$ to the connected components of $X$.
If $\mathcal{E}$ is just a coherent sheaf, there are several different
ways of extending the definition\footnote{If $X$ is a proper curve
and $\mathcal{F}$ is a coherent sheaf on $X$, then one often defines
the degree as $\chi(X, \mathcal{F}) - r\chi(X, \mathcal{O}_X)$
where $r = \dim_{\kappa(\xi)} \mathcal{F}_\xi$ is the rank of $\mathcal{F}$
at the generic point $\xi$ of $X$.}.
In a series of lemmas we show that this definition has all the properties
one expects of the degree.

\begin{lemma}
\label{lemma-degree-base-change}
Let $k'/k$ be an extension of fields. Let $X$ be a proper scheme of
dimension $\leq 1$ over $k$. Let $\mathcal{E}$ be a locally free
$\mathcal{O}_X$-module of constant rank $n$. Then the degree of
$\mathcal{E}/X/k$ is equal to the degree of
$\mathcal{E}_{k'}/X_{k'}/k'$.
\end{lemma}

\begin{proof}
More precisely, set $X_{k'} = X \times_{\Spec(k)} \Spec(k')$.
Let $\mathcal{E}_{k'} = p^*\mathcal{E}$ where $p : X_{k'} \to X$
is the projection. By
Cohomology of Schemes, Lemma \ref{coherent-lemma-flat-base-change-cohomology}
we have
$H^i(X_{k'}, \mathcal{E}_{k'}) = H^i(X, \mathcal{E}) \otimes_k k'$
and
$H^i(X_{k'}, \mathcal{O}_{X_{k'}}) = H^i(X, \mathcal{O}_X) \otimes_k k'$.
Hence we see that the Euler characteristics are unchanged, hence the
degree is unchanged.
\end{proof}

\begin{lemma}
\label{lemma-degree-additive}
Let $k$ be a field. Let $X$ be a proper scheme of dimension $\leq 1$
over $k$. Let $0 \to \mathcal{E}_1 \to \mathcal{E}_2 \to \mathcal{E}_3 \to 0$
be a short exact sequence of locally free $\mathcal{O}_X$-modules
each of finite constant rank. Then
$$
\deg(\mathcal{E}_2) = \deg(\mathcal{E}_1) + \deg(\mathcal{E}_3)
$$
\end{lemma}

\begin{proof}
Follows immediately from additivity of Euler characteristics
(Lemma \ref{lemma-euler-characteristic-additive})
and additivity of ranks.
\end{proof}

\begin{lemma}
\label{lemma-degree-birational-pullback}
Let $k$ be a field. Let $f : X' \to X$ be a birational morphism of
proper schemes of dimension $\leq 1$ over $k$. Then
$$
\deg(f^*\mathcal{E}) = \deg(\mathcal{E})
$$
for every finite locally free sheaf of constant rank. More generally
it suffices if $f$ induces a bijection between irreducible components
of dimension $1$ and isomorphisms of local rings at the corresponding
generic points.
\end{lemma}

\begin{proof}
The morphism $f$ is proper
(Morphisms, Lemma \ref{morphisms-lemma-image-proper-scheme-closed})
and has fibres of dimension $\leq 0$. Hence $f$ is finite
(Cohomology of Schemes, Lemma
\ref{coherent-lemma-proper-finite-fibre-finite-in-neighbourhood}).
Thus
$$
Rf_*f^*\mathcal{E} = f_*f^*\mathcal{E} =
\mathcal{E} \otimes_{\mathcal{O}_X} f_*\mathcal{O}_{X'}
$$
Since $f$ induces an isomorphism on local rings at generic points of
all irreducible components of dimension $1$ we see that the kernel
and cokernel
$$
0 \to \mathcal{K} \to \mathcal{O}_X \to f_*\mathcal{O}_{X'}
\to \mathcal{Q} \to 0
$$
have supports of dimension $\leq 0$. Note that tensoring this with
$\mathcal{E}$ is still an exact sequence as $\mathcal{E}$ is locally free.
We obtain
\begin{align*}
\chi(X, \mathcal{E}) - \chi(X', f^*\mathcal{E})
& =
\chi(X, \mathcal{E}) - \chi(X, f_*f^*\mathcal{E}) \\
& =
\chi(X, \mathcal{E}) - \chi(X, \mathcal{E} \otimes f_*\mathcal{O}_{X'}) \\
& =
\chi(X, \mathcal{K} \otimes \mathcal{E}) -
\chi(X, \mathcal{Q} \otimes \mathcal{E}) \\
& =
n\chi(X, \mathcal{K}) -
n\chi(X, \mathcal{Q}) \\
& =
n\chi(X, \mathcal{O}_X) - n\chi(X, f_*\mathcal{O}_{X'}) \\
& =
n\chi(X, \mathcal{O}_X) - n\chi(X', \mathcal{O}_{X'})
\end{align*}
which proves what we want. The first equality as $f$ is finite, see
Cohomology of Schemes, Lemma \ref{coherent-lemma-relative-affine-cohomology}.
The second equality by projection formula, see
Cohomology, Lemma \ref{cohomology-lemma-projection-formula}.
The third by additivity of Euler characteristics, see
Lemma \ref{lemma-euler-characteristic-additive}.
The fourth by Lemma \ref{lemma-chi-tensor-finite}.
\end{proof}

\begin{lemma}
\label{lemma-degree-on-proper-curve}
Let $k$ be a field. Let $X$ be a proper curve over $k$ with generic point
$\xi$. Let $\mathcal{E}$ be a locally free $\mathcal{O}_X$-module of rank $n$
and let $\mathcal{F}$ be a coherent $\mathcal{O}_X$-module. Then
$$
\chi(X, \mathcal{E} \otimes \mathcal{F}) =
r \deg(\mathcal{E}) + n \chi(X, \mathcal{F})
$$
where $r = \dim_{\kappa(\xi)} \mathcal{F}_\xi$ is the rank of $\mathcal{F}$.
\end{lemma}

\begin{proof}
Let $\mathcal{P}$ be the property of coherent sheaves $\mathcal{F}$
on $X$ expressing that the formula of the lemma holds. We claim that
the assumptions (1) and (2) of
Cohomology of Schemes, Lemma \ref{coherent-lemma-property}
hold for $\mathcal{P}$. Namely, (1) holds because the Euler characteristic
and the rank $r$ are additive in short exact sequences of coherent sheaves.
And (2) holds too: If $Z = X$ then we may take
$\mathcal{G} = \mathcal{O}_X$ and $\mathcal{P}(\mathcal{O}_X)$
is true by the definition of degree. If $i : Z \to X$ is the inclusion
of a closed point we may take $\mathcal{G} = i_*\mathcal{O}_Z$
and $\mathcal{P}$ holds by Lemma \ref{lemma-chi-tensor-finite}
and the fact that $r = 0$ in this case.
\end{proof}

\noindent
Let $k$ be a field. Let $X$ be a finite type scheme over $k$ of
dimension $\leq 1$. Let $C_i \subset X$, $i = 1, \ldots, t$ be
the irreducible components of dimension $1$. We view $C_i$ as a scheme
by using the induced reduced scheme structure. Let $\xi_i \in C_i$ be the
generic point. The {\it multiplicity of $C_i$ in $X$} is defined
as the length
$$
m_i = \text{length}_{\mathcal{O}_{X, \xi_i}} \mathcal{O}_{X, \xi_i}
$$
This makes sense because $\mathcal{O}_{X, \xi_i}$ is a zero dimensional
Noetherian local ring and hence has finite length over itself
(Algebra, Proposition \ref{algebra-proposition-dimension-zero-ring}).
See Chow Homology, Section \ref{chow-section-cycle-of-closed-subscheme}
for additional information. It turns out the degree of a locally free
sheaf only depends on the restriction of the irreducible components.

\begin{lemma}
\label{lemma-degree-in-terms-of-components}
Let $k$ be a field. Let $X$ be a proper scheme of dimension $\leq 1$ over $k$.
Let $\mathcal{E}$ be a locally free $\mathcal{O}_X$-module of rank $n$.
Then
$$
\deg(\mathcal{E}) = \sum m_i \deg(\mathcal{E}|_{C_i})
$$
where $C_i \subset X$, $i = 1, \ldots, t$ are the irreducible components
of dimension $1$ with reduced induced scheme structure and $m_i$ is the
multiplicity of $C_i$ in $X$.
\end{lemma}

\begin{proof}
Observe that the statement makes sense because $C_i \to \Spec(k)$
is proper of dimension $1$; we will use below that any closed subscheme
of $X$ is proper over $k$ by Morphisms, Lemmas
\ref{morphisms-lemma-closed-immersion-proper} and
\ref{morphisms-lemma-composition-proper}. Consider the open subscheme
$U_i = X \setminus (\bigcup_{j \not = i} C_j)$ and let $X_i \subset X$
be the scheme theoretic closure of $U_i$. Note that $X_i \cap U_i = U_i$
(scheme theoretically) and that $X_i \cap U_j = \emptyset$
(set theoretically) for $i \not = j$; this follows from the description
of scheme theoretic closure in
Morphisms, Lemma \ref{morphisms-lemma-quasi-compact-immersion}.
Thus we may apply Lemma \ref{lemma-degree-birational-pullback} to the morphism
$X' = \coprod X_i \to X$. Since it is clear that $C_i \subset X_i$
(scheme theoretically) and that the multiplicity of $C_i$ in $X_i$ is equal
to the multiplicity of $C_i$ in $X$, we see that we reduce to the case
discussed in the following paragraph.

\medskip\noindent
Assume $X$ is irreducible with generic point $\xi$. Let
$C = X_{red}$ have multiplicity $m$.
We have to show that $\deg(\mathcal{E}) = m \deg(\mathcal{E}|_C)$.
Let $\mathcal{I} \subset \mathcal{O}_X$ be the ideal defining the closed
subscheme $C$. Let $e \geq 0$ be minimal such that $\mathcal{I}^{e + 1} = 0$
(Cohomology of Schemes, Lemma \ref{coherent-lemma-power-ideal-kills-sheaf}).
We argue by induction on $e$. If $e = 0$, then $X = C$ and the result
is immediate. Otherwise we set $\mathcal{F} = \mathcal{I}^e$ viewed as
a coherent $\mathcal{O}_C$-module (Cohomology of Schemes, Lemma
\ref{coherent-lemma-i-star-equivalence}).
Let $X' \subset X$ be the closed subscheme cut out by the
coherent ideal $\mathcal{I}^e$ and let $m'$ be the multiplicity
of $C$ in $X'$. Taking stalks at $\xi$ of the short exact sequence
$$
0 \to \mathcal{F} \to \mathcal{O}_X \to \mathcal{O}_{X'} \to 0
$$
we find (use Algebra, Lemmas \ref{algebra-lemma-length-additive},
\ref{algebra-lemma-dimension-is-length}, and
\ref{algebra-lemma-length-independent}) that
$$
m = \text{length}_{\mathcal{O}_{X, \xi}} \mathcal{O}_{X, \xi}
= \dim_{\kappa(\xi)} \mathcal{F}_\xi +
\text{length}_{\mathcal{O}_{X', \xi}} \mathcal{O}_{X', \xi}
= r + m'
$$
where $r$ is the rank of $\mathcal{F}$ as a coherent sheaf on $C$.
Tensoring with $\mathcal{E}$ we obtain a short exact sequence
$$
0 \to \mathcal{E}|_C \otimes \mathcal{F} \to \mathcal{E} \to
\mathcal{E} \otimes \mathcal{O}_{X'} \to 0
$$
By induction we have
$\deg(\mathcal{E}|_{X'}) = m' \deg(\mathcal{E}|_C)$.
By Lemma \ref{lemma-degree-on-proper-curve} we have
$\chi(\mathcal{E}|_C \otimes \mathcal{F}) =
r \deg(\mathcal{E}|_C) + n \chi(\mathcal{F})$.
Putting everything together we obtain the result.
\end{proof}

\begin{lemma}
\label{lemma-degree-tensor-product}
Let $k$ be a field, let $X$ be a proper scheme of dimension $\leq 1$
over $k$, and let $\mathcal{E}$, $\mathcal{V}$ be locally free
$\mathcal{O}_X$-modules of constant finite rank. Then
$$
\deg(\mathcal{E} \otimes \mathcal{V}) =
\text{rank}(\mathcal{E}) \deg(\mathcal{V}) +
\text{rank}(\mathcal{V}) \deg(\mathcal{E})
$$
\end{lemma}

\begin{proof}
By Lemma \ref{lemma-degree-in-terms-of-components} and elementary
arithmetic, we reduce to the case of a proper curve.
This case follows from Lemma \ref{lemma-degree-on-proper-curve}.
\end{proof}

\begin{lemma}
\label{lemma-degree-and-det}
Let $k$ be a field, let $X$ be a proper scheme of dimension $\leq 1$
over $k$, and let $\mathcal{E}$ be a locally free
$\mathcal{O}_X$-module of rank $n$. Then
$$
\deg(\mathcal{E}) = \deg(\wedge^n(\mathcal{E})) = \deg(\det(\mathcal{E}))
$$
\end{lemma}

\begin{proof}
By Lemma \ref{lemma-degree-in-terms-of-components} and elementary
arithmetic, we reduce to the case of a proper curve.
Then there exists a modification $f : X' \to X$ such that
$f^*\mathcal{E}$ has a filtration whose successive
quotients are invertible modules, see
Divisors, Lemma \ref{divisors-lemma-filter-after-modification}.
By Lemma \ref{lemma-degree-birational-pullback} we may work on $X'$.
Thus we may assume we have a filtration
$$
0 = \mathcal{E}_0 \subset \mathcal{E}_1 \subset \mathcal{E}_2 \subset
\ldots \subset \mathcal{E}_n = \mathcal{E}
$$
by locally free $\mathcal{O}_X$-modules
with $\mathcal{L}_i = \mathcal{E}_i/\mathcal{E}_{i - 1}$ is invertible.
By Modules, Lemma \ref{modules-lemma-det-ses} and induction we find
$\det(\mathcal{E}) = \mathcal{L}_1 \otimes \ldots \otimes \mathcal{L}_n$.
Thus the equality follows from Lemma \ref{lemma-degree-tensor-product}
and additivity (Lemma \ref{lemma-degree-additive}).
\end{proof}

\begin{lemma}
\label{lemma-degree-effective-Cartier-divisor}
Let $k$ be a field, let $X$ be a proper scheme of dimension $\leq 1$
over $k$. Let $D$ be an effective Cartier divisor on $X$.
Then $D$ is finite over $\Spec(k)$ of degree
$\deg(D) = \dim_k \Gamma(D, \mathcal{O}_D)$. For a locally free sheaf
$\mathcal{E}$ of rank $n$ we have
$$
\deg(\mathcal{E}(D)) = n\deg(D) + \deg(\mathcal{E})
$$
where $\mathcal{E}(D) = \mathcal{E} \otimes_{\mathcal{O}_X} \mathcal{O}_X(D)$.
\end{lemma}

\begin{proof}
Since $D$ is nowhere dense in $X$ (Divisors, Lemma
\ref{divisors-lemma-complement-effective-Cartier-divisor})
we see that $\dim(D) \leq 0$. Hence $D$ is finite over $k$
by Lemma \ref{lemma-algebraic-scheme-dim-0}. Since $k$ is a field,
the morphism $D \to \Spec(k)$ is finite locally free and hence has
a degree
(Morphisms, Definition \ref{morphisms-definition-finite-locally-free}),
which is clearly equal to $\dim_k \Gamma(D, \mathcal{O}_D)$
as stated in the lemma. By
Divisors, Definition
\ref{divisors-definition-invertible-sheaf-effective-Cartier-divisor}
there is a short exact sequence
$$
0 \to \mathcal{O}_X \to \mathcal{O}_X(D) \to i_*i^*\mathcal{O}_X(D) \to 0
$$
where $i : D \to X$ is the closed immersion. Tensoring with $\mathcal{E}$
we obtain a short exact sequence
$$
0 \to \mathcal{E} \to \mathcal{E}(D) \to i_*i^*\mathcal{E}(D) \to 0
$$
The equation of the lemma follows from additivity of
the Euler characteristic (Lemma \ref{lemma-euler-characteristic-additive})
and Lemma \ref{lemma-chi-tensor-finite}.
\end{proof}

\begin{lemma}
\label{lemma-divisible}
Let $k$ be a field. Let $X$ be a proper scheme over $k$
which is reduced and connected.
Let $\kappa = H^0(X, \mathcal{O}_X)$. Then $\kappa/k$ is a
finite extension of fields and $w = [\kappa : k]$ divides
\begin{enumerate}
\item $\deg(\mathcal{E})$ for all locally free $\mathcal{O}_X$-modules
$\mathcal{E}$,
\item $[\kappa(x) : k]$ for all closed points $x \in X$, and
\item $\deg(D)$ for all closed subschemes $D \subset X$
of dimension zero.
\end{enumerate}
\end{lemma}

\begin{proof}
See Lemma \ref{lemma-proper-geometrically-reduced-global-sections}
for the assertions about $\kappa$.
For every quasi-coherent $\mathcal{O}_X$-module, the
$k$-vector spaces $H^i(X, \mathcal{F})$ are $\kappa$-vector spaces.
The divisibilities easily follow from this statement and the
definitions.
\end{proof}

\begin{lemma}
\label{lemma-degree-pullback-map-proper-curves}
Let $k$ be a field. Let $f : X \to Y$ be a nonconstant morphism of
proper curves over $k$. Let $\mathcal{E}$ be a locally free
$\mathcal{O}_Y$-module. Then
$$
\deg(f^*\mathcal{E}) = \deg(X/Y) \deg(\mathcal{E})
$$
\end{lemma}

\begin{proof}
The degree of $X$ over $Y$ is defined in
Morphisms, Definition \ref{morphisms-definition-degree}.
Thus $f_*\mathcal{O}_X$ is a coherent $\mathcal{O}_Y$-module
of rank $\deg(X/Y)$, i.e.,
$\deg(X/Y) = \dim_{\kappa(\xi)} (f_*\mathcal{O}_X)_\xi$ where $\xi$
is the generic point of $Y$. Thus we obtain
\begin{align*}
\chi(X, f^*\mathcal{E})
& =
\chi(Y, f_*f^*\mathcal{E}) \\
& =
\chi(Y, \mathcal{E} \otimes f_*\mathcal{O}_X) \\
& =
\deg(X/Y) \deg(\mathcal{E}) + n \chi(Y, f_*\mathcal{O}_X) \\
& =
\deg(X/Y) \deg(\mathcal{E}) + n \chi(X, \mathcal{O}_X)
\end{align*}
as desired. The first equality as $f$ is finite, see
Cohomology of Schemes, Lemma \ref{coherent-lemma-relative-affine-cohomology}.
The second equality by projection formula, see
Cohomology, Lemma \ref{cohomology-lemma-projection-formula}.
The third equality by Lemma \ref{lemma-degree-on-proper-curve}.
\end{proof}

\noindent
The following is a trivial but important consequence of the
results on degrees above.

\begin{lemma}
\label{lemma-check-invertible-sheaf-trivial}
Let $k$ be a field. Let $X$ be a proper curve over $k$.
Let $\mathcal{L}$ be an invertible $\mathcal{O}_X$-module.
\begin{enumerate}
\item If $\mathcal{L}$ has a nonzero section, then
$\deg(\mathcal{L}) \geq 0$.
\item If $\mathcal{L}$ has a nonzero section $s$ which vanishes
at a point, then $\deg(\mathcal{L}) > 0$.
\item If $\mathcal{L}$ and $\mathcal{L}^{-1}$ have nonzero sections, then
$\mathcal{L} \cong \mathcal{O}_X$.
\item If $\deg(\mathcal{L}) \leq 0$ and $\mathcal{L}$ has a nonzero
section, then $\mathcal{L} \cong \mathcal{O}_X$.
\item If $\mathcal{N} \to \mathcal{L}$ is a nonzero map of invertible
$\mathcal{O}_X$-modules, then $\deg(\mathcal{L}) \geq \deg(\mathcal{N})$
and if equality holds then it is an isomorphism.
\end{enumerate}
\end{lemma}

\begin{proof}
Let $s$ be a nonzero section of $\mathcal{L}$. Since $X$ is a curve, we
see that $s$ is a regular section. Hence there is an effective
Cartier divisor $D \subset X$ and an isomorphism
$\mathcal{L} \to \mathcal{O}_X(D)$ mapping $s$ the canonical
section $1$ of $\mathcal{O}_X(D)$, see
Divisors, Lemma \ref{divisors-lemma-characterize-OD}.
Then $\deg(\mathcal{L}) = \deg(D)$ by
Lemma \ref{lemma-degree-effective-Cartier-divisor}.
As $\deg(D) \geq 0$ and $= 0$
if and only if $D = \emptyset$, this proves (1) and (2).
In case (3) we see that $\deg(\mathcal{L}) = 0$ and
$D = \emptyset$. Similarly for (4). To see (5) apply
(1) and (4) to the invertible sheaf
$$
\mathcal{L} \otimes_{\mathcal{O}_X} \mathcal{N}^{\otimes -1} =
\SheafHom_{\mathcal{O}_X}(\mathcal{N}, \mathcal{L})
$$
which has degree $\deg(\mathcal{L}) - \deg(\mathcal{N})$
by Lemma \ref{lemma-degree-tensor-product}.
\end{proof}

\begin{lemma}
\label{lemma-no-sections-dual-nef}
Let $k$ be a field. Let $X$ be a proper scheme over $k$
which is reduced, connected, and equidimensional of dimension $1$.
Let $\mathcal{L}$ be an invertible $\mathcal{O}_X$-module.
If $\deg(\mathcal{L}|_C) \leq 0$ for all irreducible components $C$
of $X$, then either $H^0(X, \mathcal{L}) = 0$ or
$\mathcal{L} \cong \mathcal{O}_X$.
\end{lemma}

\begin{proof}
Let $s \in H^0(X, \mathcal{L})$ be nonzero. Since $X$ is reduced there exists
an irreducible component $C$ of $X$ with $s|_C \not = 0$.
But if $s|_C$ is nonzero, then
$s$ is nonwhere vanishing on $C$ by
Lemma \ref{lemma-check-invertible-sheaf-trivial}.
This in turn implies $s$ is nowhere vanishing on
every irreducible component of $X$ meeting $C$.
Since $X$ is connected, we conclude that $s$
vanishes nowhere and the lemma follows.
\end{proof}

\begin{lemma}
\label{lemma-ample-curve}
Let $k$ be a field. Let $X$ be a proper curve over $k$.
Let $\mathcal{L}$ be an invertible $\mathcal{O}_X$-module.
Then $\mathcal{L}$ is ample if and only if $\deg(\mathcal{L}) > 0$.
\end{lemma}

\begin{proof}
If $\mathcal{L}$ is ample, then there exists an $n > 0$ and a section
$s \in H^0(X, \mathcal{L}^{\otimes n})$ with $X_s$ affine. Since
$X$ isn't affine (otherwise by
Morphisms, Lemma \ref{morphisms-lemma-finite-proper}
$X$ would be finite), we see that $s$ vanishes at some point.
Hence $\deg(\mathcal{L}^{\otimes n}) > 0$ by
Lemma \ref{lemma-check-invertible-sheaf-trivial}.
By Lemma \ref{lemma-degree-tensor-product}
we conclude that $\deg(\mathcal{L}) = 1/n\deg(\mathcal{L}^{\otimes n}) > 0$.

\medskip\noindent
Assume $\deg(\mathcal{L}) > 0$. Then
$$
\dim_k H^0(X, \mathcal{L}^{\otimes n}) \geq \chi(X, \mathcal{L}^n)
= n\deg(\mathcal{L}) + \chi(X, \mathcal{O}_X)
$$
grows linearly with $n$. Hence for any finite collection of closed
points $x_1, \ldots, x_t$ of $X$, we can find an $n$ such that
$\dim_k H^0(X, \mathcal{L}^{\otimes n}) > \sum \dim_k \kappa(x_i)$.
(Recall that by Hilbert Nullstellensatz, the extension fields
$\kappa(x_i)/k$ are finite, see for example
Morphisms, Lemma \ref{morphisms-lemma-closed-point-fibre-locally-finite-type}).
Hence we can find a nonzero $s \in H^0(X, \mathcal{L}^{\otimes n})$
vanishing in $x_1, \ldots, x_t$. In particular, if we choose
$x_1, \ldots, x_t$ such that $X \setminus \{x_1, \ldots, x_t\}$
is affine, then $X_s$ is affine too
(for example by Properties, Lemma \ref{properties-lemma-affine-cap-s-open}
although if we choose our finite set such that
$\mathcal{L}|_{X \setminus \{x_1, \ldots, x_t\}}$ is trivial, then
it is immediate). The conclusion is that we can find an $n > 0$
and a nonzero section $s \in H^0(X, \mathcal{L}^{\otimes n})$ such
that $X_s$ is affine.

\medskip\noindent
We will show that for every quasi-coherent sheaf of ideals $\mathcal{I}$
there exists an $m > 0$ such that
$H^1(X, \mathcal{I} \otimes \mathcal{L}^{\otimes m})$ is zero.
This will finish the proof by
Cohomology of Schemes, Lemma \ref{coherent-lemma-vanshing-gives-ample}.
To see this we consider the maps
$$
\mathcal{I} \xrightarrow{s}
\mathcal{I} \otimes \mathcal{L}^{\otimes n} \xrightarrow{s}
\mathcal{I} \otimes \mathcal{L}^{\otimes 2n} \xrightarrow{s} \ldots
$$
Since $\mathcal{I}$ is torsion free, these maps are injective and
isomorphisms over $X_s$, hence the cokernels have vanishing $H^1$
(by Cohomology of Schemes, Lemma
\ref{coherent-lemma-coherent-support-dimension-0} for example).
We conclude that the maps of vector spaces
$$
H^1(X, \mathcal{I}) \to
H^1(X, \mathcal{I} \otimes \mathcal{L}^{\otimes n}) \to
H^1(X, \mathcal{I} \otimes \mathcal{L}^{\otimes 2n}) \to \ldots
$$
are surjective. On the other hand, the dimension of $H^1(X, \mathcal{I})$
is finite, and every element maps to zero eventually by
Cohomology of Schemes, Lemma
\ref{coherent-lemma-section-affine-open-kills-classes}.
Thus for some $e > 0$ we see that
$H^1(X, \mathcal{I} \otimes \mathcal{L}^{\otimes en})$ is zero.
This finishes the proof.
\end{proof}

\begin{lemma}
\label{lemma-ampleness-in-terms-of-degrees-components}
Let $k$ be a field. Let $X$ be a proper scheme of dimension $\leq 1$
over $k$. Let $\mathcal{L}$ be an invertible $\mathcal{O}_X$-module.
Let $C_i \subset X$, $i = 1, \ldots, t$ be the irreducible components
of dimension $1$. The following are equivalent:
\begin{enumerate}
\item $\mathcal{L}$ is ample, and
\item $\deg(\mathcal{L}|_{C_i}) > 0$ for $i = 1, \ldots, t$.
\end{enumerate}
\end{lemma}

\begin{proof}
Let $x_1, \ldots, x_r \in X$ be the isolated closed points.
Think of $x_i = \Spec(\kappa(x_i))$ as a scheme.
Consider the morphism of schemes
$$
f : C_1 \amalg \ldots \amalg C_t \amalg x_1 \amalg \ldots \amalg x_r
\longrightarrow X
$$
This is a finite surjective morphism of schemes proper over $k$
(details omitted). Thus $\mathcal{L}$ is ample if and only if
$f^*\mathcal{L}$ is ample (Cohomology of Schemes, Lemma
\ref{coherent-lemma-surjective-finite-morphism-ample}).
Thus we conclude by Lemma \ref{lemma-ample-curve}.
\end{proof}

\begin{lemma}
\label{lemma-general-degree-g-line-bundle}
Let $k$ be an algebraically closed field. Let $X$ be a proper curve over $k$.
Then there exist
\begin{enumerate}
\item an invertible $\mathcal{O}_X$-module $\mathcal{L}$ with
$\dim_k H^0(X, \mathcal{L}) = 1$ and $H^1(X, \mathcal{L}) = 0$, and
\item an invertible $\mathcal{O}_X$-module $\mathcal{N}$ with
$\dim_k H^0(X, \mathcal{N}) = 0$ and $H^1(X, \mathcal{N}) = 0$.
\end{enumerate}
\end{lemma}

\begin{proof}
Choose a closed immersion $i : X \to \mathbf{P}^n_k$
(Lemma \ref{lemma-dim-1-proper-projective}).
Setting $\mathcal{L} = i^*\mathcal{O}_{\mathbf{P}^n}(d)$
for $d \gg 0$ we see that there exists an invertible sheaf
$\mathcal{L}$ with $H^0(X, \mathcal{L}) \not = 0$ and
$H^1(X, \mathcal{L}) = 0$ (see
Cohomology of Schemes, Lemma \ref{coherent-lemma-vanshing-gives-ample}
for vanishing and the references therein for nonvanishing).
We will finish the proof of (1) by descending induction on
$t = \dim_k H^0(X, \mathcal{L})$. The base case $t = 1$ is trivial.
Assume $t > 1$.

\medskip\noindent
Let $U \subset X$ be the nonempty open subset of nonsingular points
studied in Lemma \ref{lemma-dense-smooth-open-variety-over-perfect-field}.
Let $s \in H^0(X, \mathcal{L})$ be nonzero. There exists a closed
point $x \in U$ such that $s$ does not vanish in $x$. Let $\mathcal{I}$
be the ideal sheaf of $i : x \to X$ as in
Lemma \ref{lemma-regular-point-on-curve}. Look at the
short exact sequence
$$
0 \to \mathcal{I} \otimes_{\mathcal{O}_X} \mathcal{L} \to
\mathcal{L} \to i_*i^*\mathcal{L} \to 0
$$
Observe that $H^0(X, i_*i^*\mathcal{L}) = H^0(x, i^*\mathcal{L})$
has dimension $1$ as $x$ is a $k$-rational point ($k$ is algebraically
closed). Since $s$ does not vanish at $x$ we conclude that
$$
H^0(X, \mathcal{L}) \longrightarrow H^0(X, i_*i^*\mathcal{L})
$$
is surjective. Hence
$\dim_k H^0(X, \mathcal{I} \otimes_{\mathcal{O}_X} \mathcal{L}) = t - 1$.
Finally, the long exact sequence of cohomology also shows that
$H^1(X, \mathcal{I} \otimes_{\mathcal{O}_X} \mathcal{L}) = 0$
thereby finishing the proof of the induction step.

\medskip\noindent
To get an invertible sheaf as in (2) take an invertible sheaf
$\mathcal{L}$ as in (1) and do the argument in the previous paragraph
one more time.
\end{proof}

\begin{lemma}
\label{lemma-vanishing-degree-2g-and-1-line-bundle}
Let $k$ be an algebraically closed field. Let $X$ be a proper curve over $k$.
Set $g = \dim_k H^1(X, \mathcal{O}_X)$. For every invertible
$\mathcal{O}_X$-module $\mathcal{L}$ with $\deg(\mathcal{L}) \geq 2g - 1$
we have $H^1(X, \mathcal{L}) = 0$.
\end{lemma}

\begin{proof}
Let $\mathcal{N}$ be the invertible module we found in
Lemma \ref{lemma-general-degree-g-line-bundle} part (2).
The degree of $\mathcal{N}$ is
$\chi(X, \mathcal{N}) - \chi(X, \mathcal{O}_X) = 0 - (1 - g) = g - 1$.
Hence the degree of $\mathcal{L} \otimes \mathcal{N}^{\otimes - 1}$
is $\deg(\mathcal{L}) - (g - 1) \geq g$.
Hence
$\chi(X, \mathcal{L} \otimes \mathcal{N}^{\otimes -1}) \geq g + 1 - g = 1$.
Thus there is a nonzero global section $s$ whose zero scheme is an
effective Cartier divisor $D$ of degree $\deg(\mathcal{L}) - (g - 1)$.
This gives a short exact sequence
$$
0 \to \mathcal{N} \xrightarrow{s} \mathcal{L} \to i_*(\mathcal{L}|_D) \to 0
$$
where $i : D \to X$ is the inclusion morphism. We conclude that
$H^0(X, \mathcal{L})$ maps isomorphically to $H^0(D, \mathcal{L}|_D)$
which has dimension $\deg(\mathcal{L}) - (g - 1)$. The result follows
from the definition of degree.
\end{proof}






